Although the present work is simulation-based, it raises practical issues that remain relevant for real multi-robot deployments.
The primary ethical concern is safety and reliability: generated decisions may be infeasible, unstable, or unproductive if they are not constrained by system-side checks.
The studied control loop executes only grounded action identifiers after candidate shaping, feasibility validation, and fallback handling.
Unsafe or clearly invalid actions are filtered before execution.
The logging of prompts, parsed outputs, executed actions, and observed step consequences also improves traceability, which is important when analyzing failures such as long no-delivery tails, stale waiting loops, or battery-related breakdowns.

From a sustainability perspective, the main trade-off is between model capability and computational cost.
Larger language models generally require more inference time and compute, while smaller models may require stronger prompt structure and validation to remain effective.
This thesis therefore treats call efficiency, unnecessary replanning, and coordination overhead as experimental outcomes with practical relevance.
Even with safeguards, limited simulator initializations, limited scenario coverage, and the absence of physical deployment still bound the scope of the reported results.
